\begin{mainquestion}{seriesmosbias}{This Series PMOS}

\questionpoints
Two identical PMOS transistors $Q_1$ and $Q_2$ with $k_p = 1\text{ mA/V}^2$ and $V_{tp} = -1\text{ V}$ are connected in series as shown in the figure below.

\begin{center}
\begin{circuitikz}[american]
    % VDD (+20V)
    \draw (0,6) node[vcc] {+20 V}
        to[R, l=$6\text{ k}\Omega$] (0,4)
        node[label={right:$V_3$}] {}
        node[pmos, anchor=S] (Q1) {$Q_1$};
    
    % Q1 Gate connection to +10V
    \draw (Q1.G) -- ++(-0.5,0) node[ocirc] {} node[left] {+10 V};

    % Q1 Drain to Q2 Source
    \draw (Q1.D) -- (0,2)
        node[label={right:$V_2$}] {}
        node[pmos, anchor=S] (Q2) {$Q_2$};
        
    % Q2 Gate connection to +5V
    \draw (Q2.G) -- ++(-0.5,0) node[ocirc] {} node[left] {+5 V};

    % Q2 Drain to Ground (V1)
    \draw (Q2.D) -- (0,0)
        node[label={right:$V_1$}] {}
        node[ground] {};
\end{circuitikz}
\end{center}

% =========================================================
\begin{partquestion}{findvoltages}
{What are the resulting DC bias voltages $V_1$, $V_2$, and $V_3$? Be sure to show all of your work.}

\begin{answer}
        Well, $V_1$ is immediately obvious—it is identically $0\text{ V}$ (ground)!
        
        Consider transistor $Q_2$, where $V_{G2} = 5\text{ V}$ and $V_{D2} = 0\text{ V}$. Since $|V_{tp}| = 1\text{ V}$, this device will be in saturation as long as it is not in cutoff, which would occur when $V_{SG} < |V_{tp}|$ or equivalently when $V_{S2} - V_{G2} < 1 \Rightarrow V_2 - 5 < 1 \Rightarrow V_2 < 6\text{ V}$.
        
        If $Q_2$ is in saturation, $Q_1$ cannot be in cutoff, or current would not flow through $Q_2$. Let us assume that both devices are in saturation and check for possible contradictions that arise with the analysis. Since both transistors have identical $k_p$ and $V_{tp}$, and for saturation $i_D = \frac{1}{2} k_p (V_{SG} - |V_{tp}|)^2$, it follows that both must have the same $V_{SG}$ if the currents are the same; that is, $V_{SG1} = V_{SG2}$.
        
        So we have the relationship:
        \begin{align*}
            V_3 - 10 &= V_2 - 5 \\
            V_3 &= V_2 + 5
        \end{align*}
        
        Additionally, we have the current equation:
        \[ i_D = \frac{1}{2} k_p (V_{SG2} - |V_{tp}|)^2 = \frac{1}{2}(1)(V_2 - 5 - 1)^2 = 0.5(V_2 - 6)^2 \]
        
        This current must be equal to the current flowing through the top resistor:
        \[ i_D = \frac{20 - V_3}{6\text{ k}\Omega} \]
        
        Equating the two expressions and substituting $V_3 = V_2 + 5$:
        \begin{align*}
            \frac{20 - (V_2 + 5)}{6} &= 0.5(V_2 - 6)^2 \\
            \frac{15 - V_2}{6} &= 0.5(V_2^2 - 12V_2 + 36) \\
            15 - V_2 &= 3(V_2^2 - 12V_2 + 36) \\
            15 - V_2 &= 3V_2^2 - 36V_2 + 108
        \end{align*}
        
        Rearranging gives the quadratic equation:
        \[ 3V_2^2 - 35V_2 + 93 = 0 \]
        
        Which has two potential solutions: $V_2 \approx 4.093\text{ V}$ and $V_2 \approx 7.573\text{ V}$.
        
        Note the first solution $4.093 < 6$, which we earlier showed is the threshold for $Q_2$ to be in saturation, so this solution is invalidated (it implies cutoff). The second solution $V_2 = 7.573\text{ V}$ is greater than $6\text{ V}$, so it leads to a consistent result.
        
        It follows directly that:
        \[ V_3 = V_2 + 5 = 7.573 + 5 = 12.573\text{ V} \]
        
        Final Voltages: $V_1 = 0\text{ V}$, $V_2 = 7.573\text{ V}$, $V_3 = 12.573\text{ V}$.
\end{answer}

\end{partquestion}

% =========================================================
\begin{partquestion}{findcurrent}
{At these bias levels, what is the current $i_D$?}

\begin{answer}
        We can use the equation for current in saturation for either $Q_1$ or $Q_2$. Let's choose $Q_1$:
        \[ i_D = \frac{1}{2} k_p (V_{SG1} - |V_{tp}|)^2 \]
        \[ i_D = 0.5(1)(V_3 - 10 - 1)^2 = 0.5(12.573 - 11)^2 \]
        \[ i_D \approx 1.24\text{ mA} \]
\end{answer}

\end{partquestion}

\end{mainquestion}
